\origpage{240}{ROBERTSON AND SEYMOUR}
We assume then that there is no such $O$-arc $F$. From this and the fact that $\alpha(G) \ge 9$ it follows that for every loop $e$ of $G$, $\bar{e}$ surrounds some cuff. Since $\alpha(G) \ne 0$ it follows that $G$ has an edge which is not a loop, $e_1$ say. If $e_1 \notin \Delta_{r-1}$ then $e_1$ satisfies the theorem, and we assume that $e_1 \subseteq \Delta_{r-1}$. Let $C$ be a cuff with $\mathscr{A}_{r-1}(C) \ne \emptyset$ and $e_1 \subseteq \Delta_{r-1}(C)$. Let $A_{r-1}(C)$ be defined as at the start of this section. If $A_{r-1}(C) \subseteq bd(\Sigma)$ then $A_{r-1}(C)$ is a cuff $C' \ne C$, since $e_1 \subseteq \Delta_{r-1}(C)$, and so $\Sigma \cong \Sigma(0, 0, 2)$, a contradiction. Thus $A_{r-1}(C) \nsubseteq bd(\Sigma)$, and hence there is an edge $e$ of $G$ with $e \subseteq A_{r-1}(C)$. We claim that $e \notin \Delta_{r-1}$. For certainly $e \notin \Delta_{r-1}(C)$; we suppose $e \subseteq \Delta_{r-1}(C')$ for some cuff $C' \ne C$. Then $A_{r-1}(C) \cap \Delta_{r-1}(C') \ne \emptyset$, but $A_{r-1}(C) \nsubseteq \Delta_{r-1}(C')$, since $A_{r-1}(C)$ is not null-homotopic in $\Sigma + \hat{C}'$. Hence $A_{r-1}(C) \cap A_{r-1}(C') \ne \emptyset$. But $A_{r-1}(C) \subseteq \mathscr{A}_r(C)$, and $A_{r-1}(C') \subseteq \mathscr{A}_r(C')$, and so $\mathscr{A}_r(C) \cap \mathscr{A}_r(C') \ne \emptyset$, contrary to (5.3). This proves our claim that $e \notin \Delta_{r-1}$. We suppose $e$ is a loop, so that $A_{r-1}(C) = \bar{e}$. But then both regions incident with $e$ are subsets of $\mathscr{A}_1(C)$, and so $\mathscr{A}_1(C)$ is not inside $A_{r-1}(C)$, a contradiction, since $r \ge 2$. Thus $e$ is not a loop. Finally, we claim that $\bar{e} \cap bd(\Sigma) = \emptyset$. For suppose $v \in V(G)$ is an end of $e$, and $v \in C'$ for some cuff $C'$. Since $|V(G) \cap C'| \le r-1$ it follows that $\mathscr{A}_{r-1}(C')$ includes every region of $G$ incident with $v$, and so $e \subseteq \Delta_{r-1}(C') \subseteq \Delta_{r-1}$, a contradiction. Thus $\bar{e} \cap bd(\Sigma) = \emptyset$, and so $e$ satisfies the theorem.

Putting these results together, we obtain

\stmt
(7.4) \textit{Let $\Sigma \cong \Sigma(a, b, c)$, where $3a + 3b + c \ge 3$, and let $G$ be a boundary-linked proper graph in $\Sigma$ with $\alpha(G) \ge \frac{9}{2}r$, where $r \ge 2$ is an integer, and $|V(G) \cap C| \le r-1$ for every cuff $C$. Then there is an edge $e$ of $G$ which is not a loop, with $\bar{e} \cap bd(\Sigma) = \emptyset$, such that $\phi_e^{-1}(G \setminus e)$ is boundary-linked in $\Sigma_e$ and}
\begin{equation*}
\alpha(\phi_e^{-1}(G \setminus e)) \ge r.
\end{equation*}
As a first application of (7.4), we deduce the following final form of our main result. (A second application appears in Section 9.)

\stmt
(7.5) \textit{Let $\Sigma \cong \Sigma(a, b, c)$, where $3a + 3b + c \ge 3$. Let $k \ge 0$ be an integer, and let $v = v(\Sigma(a, b, c+1), k)$. Let $G$ be a boundary-linked proper graph in $\Sigma$ with $\alpha(G) \ge 42kv$ and with $|V(G) \cap bd(\Sigma)| \le 2k - 2c - 2$. Let $H$ be a forest in $\Sigma$ with $V(G) \cap bd(\Sigma) = V(H) \cap bd(\Sigma)$. Then $H$ is $G$-feasible.}

\textit{Proof.} If $c \ge 2$ this follows from (6.1), since $v(\Sigma, k) \le v$. If $V(H) \cap bd(\Sigma) = \emptyset$ the theorem is trivial. We assume then that $V(H) \cap bd(\Sigma) \ne \emptyset$ and $c \le 1$. It follows that $c = 1$ and $k \ge 3$. Let $r = \lfloor 28kv/3 \rfloor$. Then $\alpha(G) \ge 9r/2$, and we may apply (7.4). Let $e$ be an edge of $G$ satisfying the
